\documentclass[12pt, a4paper]{article}

% Paquetes
\usepackage[utf8]{inputenc}
\usepackage[T1]{fontenc}
\usepackage[spanish]{babel}
\usepackage{amsmath, amssymb}
\usepackage{geometry}
\usepackage{booktabs}
\usepackage{array}
\usepackage{hyperref}
\usepackage{xcolor}
\usepackage{enumitem}
\usepackage{fancyhdr}
\usepackage{titlesec}
\usepackage{lmodern}
\usepackage{microtype}
\usepackage{tikz}
\usepackage{pgfplots}
\pgfplotsset{compat=1.18}
\usetikzlibrary{arrows.meta, decorations.pathmorphing, backgrounds, positioning}

\geometry{top=2.5cm, bottom=2.5cm, left=3cm, right=3cm}


\definecolor{azulms}{RGB}{25, 100, 180}
\definecolor{rojoquench}{RGB}{190, 40, 40}
\definecolor{verdedisco}{RGB}{30, 140, 80}
\definecolor{naranjabulge}{RGB}{210, 110, 20}
\definecolor{moradostarburst}{RGB}{120, 50, 160}
\definecolor{grisclaro}{RGB}{245, 245, 248}


\pagestyle{fancy}
\fancyhf{}
\rhead{\textcolor{azulms}{\small Secuencia Principal de Galaxias}}
\lhead{\textcolor{azulms}{\small Astronomía Extragaláctica}}
\cfoot{\thepage}
\renewcommand{\headrulewidth}{0.5pt}


\titleformat{\section}{\large\bfseries\color{azulms}}{}{0em}{}[\titlerule]
\titleformat{\subsection}{\normalsize\bfseries\color{rojoquench}}{}{0em}{}



\begin{document}


% ── Portada
\begin{titlepage}
    \centering
    \vspace*{2cm}
    {\color{azulms}\rule{\linewidth}{1.5pt}}\\[0.4cm]
    {\LARGE\bfseries\color{azulms}
        Secuencia Principal\\[0.3cm]
        de Galaxias}\\[0.4cm]
    {\color{azulms}\rule{\linewidth}{1.5pt}}\\[1.5cm]
    {\normalsize
        Basada en Leslie et al.\ (2020)}\\[2cm]
    {\normalsize Notas de prácticas}\\[0.5cm]
    {\normalsize \today}
    \vfill
\end{titlepage}

% ── Índice ───────────────────────────────────────────────────────────────────
\tableofcontents
\newpage

% ════════════════════════════════════════════════════════════════════════════
\section{¿Qué es la Secuencia Principal?}
% ════════════════════════════════════════════════════════════════════════════

Cuando se representa la \textbf{Tasa de Formación Estelar} (SFR,
\textit{Star Formation Rate}) frente a la \textbf{masa estelar} $M_\star$
de una muestra de galaxias en escala logarítmica, se observa que las galaxias
activamente formando estrellas no se distribuyen de forma aleatoria:
la mayoría ocupa una banda bien definida y estrecha llamada la
\textbf{Secuencia Principal} (\textit{Main Sequence}, MS).

\vspace{0.5em}
\begin{center}
\begin{tikzpicture}
\begin{axis}[
    width=12cm, height=8cm,
    xlabel={$\log_{10}(M_\star\,/\,M_\odot)$},
    ylabel={$\log_{10}(\mathrm{SFR}\,/\,M_\odot\,\mathrm{yr}^{-1})$},
    xmin=8.5, xmax=12.0,
    ymin=-2.5, ymax=3.5,
    grid=both, grid style={gray!20},
    tick align=inside,
    legend pos=north west,
    legend style={font=\small},
    title={\textbf{Plano SFR -- $M_\star$}},
    title style={color=azulms}
]

%% ── Quiescent galaxies ──────────────────────────────────────────────────
\addplot[
    only marks, mark=*, mark size=1.8pt,
    color=rojoquench, opacity=0.5
] table {
9.0  -1.5
9.5  -1.2
10.0 -0.8
10.5 -0.6
11.0 -0.3
11.5  0.1
9.2  -1.8
9.8  -1.0
10.3 -0.9
10.8 -0.4
11.2 -0.1
};
\addlegendentry{Galaxias quiescentes}

%% ── Star-forming main sequence ─────────────────────────────────────────
\addplot[
    only marks, mark=square*, mark size=1.8pt,
    color=azulms, opacity=0.6
] table {
9.0   0.3
9.3   0.6
9.6   0.9
9.9   1.2
10.2  1.4
10.5  1.6
10.8  1.75
11.1  1.85
11.4  1.90
9.1   0.4
9.4   0.7
9.7   1.0
10.0  1.25
10.3  1.45
10.6  1.65
10.9  1.78
11.2  1.88
};
\addlegendentry{Secuencia Principal (SF)}

%% ── Starbursts ──────────────────────────────────────────────────────────
\addplot[
    only marks, mark=triangle*, mark size=2pt,
    color=moradostarburst, opacity=0.8
] table {
10.0  2.5
10.5  2.7
10.2  2.8
9.8   2.4
11.0  3.0
};
\addlegendentry{Starbursts}

%% ── MS fit line ─────────────────────────────────────────────────────────
\addplot[
    thick, color=azulms,
    domain=8.5:11.5, samples=100
] {min(1.92, (x - 9.0)*0.95 - 0.5)};
\addlegendentry{Ajuste MS (z $\sim$ 1)}

%% ── Quenched region ─────────────────────────────────────────────────────
\addplot[
    thick, dashed, color=rojoquench,
    domain=8.5:11.5, samples=100
] {min(0.92, (x - 9.0)*0.95 - 1.5)};
\addlegendentry{MS $-$ 1\,dex (límite quiescentes)}

%% ── Annotations ─────────────────────────────────────────────────────────
\node[font=\small\bfseries, color=azulms]
    at (axis cs: 10.5, 1.1) {$\sim$0.3\,dex de dispersión};
\draw[->, azulms, thick]
    (axis cs: 10.2, 1.05) -- (axis cs: 10.5, 1.62);

\end{axis}
\end{tikzpicture}
\end{center}

La relación muestra que, en promedio:
\begin{equation}
    \log_{10}(\mathrm{SFR}) \approx \alpha\,\log_{10}(M_\star) + \beta,
    \label{eq:ms_lineal}
\end{equation}
con pendiente $\alpha \approx 0.75$--$1.0$ y dispersión $\sigma \approx 0.2$--$0.3$\,dex.
Esta relación se mantiene observacionalmente desde $z \sim 0$ hasta al menos $z \sim 4$--$5$
(Noeske et al.\ 2007; Karim et al.\ 2011; Leslie et al.\ 2020).

% ════════════════════════════════════════════════════════════════════════════
\section{Las Dos Poblaciones de Galaxias}
% ════════════════════════════════════════════════════════════════════════════

En el plano SFR--$M_\star$ se identifican claramente dos poblaciones:

\begin{center}
\begin{tabular}{p{4cm} p{4.5cm} p{4.5cm}}
\toprule
\textbf{Propiedad} & \textbf{Star-Forming (MS)} & \textbf{Quiescentes} \\
\midrule
Posición en el plano & Sobre la MS & $\sim 1$\,dex bajo la MS \\
Color óptico & Azul & Roja \\
Morfología & Disco, espiral & Elíptica, esferical \\
Gas & Rica en gas frío & Pobre en gas \\
sSFR & Alto & Muy bajo \\
Masa típica & Cualquier masa & Alta masa \\
\bottomrule
\end{tabular}
\end{center}

Además existe una pequeña población de \textbf{Starbursts} ($\sim$5--20\% de las SF)
que se ubica $\sim$1\,dex \textit{sobre} la MS, con SFRs muy elevados tipicamente
asociados a fusiones (\textit{mergers}) de galaxias.

% ════════════════════════════════════════════════════════════════════════════
\section{La Tasa Específica de Formación Estelar (sSFR)}
% ════════════════════════════════════════════════════════════════════════════

El \textbf{sSFR} (\textit{specific Star Formation Rate}) se define como:

\begin{equation}
    \mathrm{sSFR} \equiv \frac{\mathrm{SFR}}{M_\star}
    \quad [\mathrm{yr}^{-1}],
    \label{eq:ssfr}
\end{equation}

y mide la \textit{eficiencia relativa} con que una galaxia convierte su masa
en nuevas estrellas. Es el inverso del tiempo de doblado de masa:

\begin{equation}
    t_{\mathrm{doble}} = \frac{1}{\mathrm{sSFR}}
    \approx \frac{M_\star}{\mathrm{SFR}}.
    \label{eq:tdouble}
\end{equation}

Propiedades clave del sSFR:
\begin{itemize}[leftmargin=2em]
    \item Las galaxias más masivas tienen \textbf{menor sSFR} a todo redshift
          (se ``apagan'' más eficientemente).
    \item El sSFR \textbf{aumenta fuertemente con el redshift}:
          \begin{equation}
              \mathrm{sSFR}(z) \propto (1+z)^{2.8}
              \quad \text{(Leslie et al.\ 2020)}.
          \end{equation}
    \item Esto implica que en el universo joven las galaxias formaban estrellas
          mucho más activamente que hoy.
\end{itemize}

\begin{center}
\begin{tikzpicture}
\begin{axis}[
    width=11cm, height=6.5cm,
    xlabel={Redshift $z$},
    ylabel={$\log_{10}(\mathrm{sSFR}\,/\,\mathrm{Gyr}^{-1})$},
    xmin=0, xmax=5,
    ymin=-2, ymax=1.5,
    grid=both, grid style={gray!20},
    legend pos=north west,
    legend style={font=\small},
    title={\textbf{Evolución del sSFR con el redshift}},
    title style={color=azulms}
]

%% Líneas para distintas masas (sSFR ∝ (1+z)^2.8, desplazadas por masa)
\addplot[thick, color=azulms, domain=0:5, samples=80]
    {2.8*log10(1+x) - 1.6};
\addlegendentry{$\log(M_\star/M_\odot)=10.1$}

\addplot[thick, color=verdedisco, dashed, domain=0:5, samples=80]
    {2.8*log10(1+x) - 1.9};
\addlegendentry{$\log(M_\star/M_\odot)=10.5$}

\addplot[thick, color=rojoquench, dotted, domain=0:5, samples=80]
    {2.8*log10(1+x) - 2.3};
\addlegendentry{$\log(M_\star/M_\odot)=11.0$}

\node[font=\small, color=azulms]
    at (axis cs: 4.0, 0.8) {$\mathrm{sSFR}\propto(1+z)^{2.8}$};

\end{axis}
\end{tikzpicture}
\end{center}

% ════════════════════════════════════════════════════════════════════════════
\section{Forma de la Secuencia Principal}
% ════════════════════════════════════════════════════════════════════════════

\subsection{lineal o con aplanamiento}

\begin{description}[leftmargin=2em]
    \item[\textcolor{azulms}{Relación lineal}]
        $\log\,\mathrm{SFR} = \alpha\,\log\,M_\star + \beta$ con $\alpha \approx 0.75$--$1.0$
        en todo el rango de masas. Típicamente encontrada en estudios que seleccionan
        galaxias muy azules o solo discos (Speagle et al.\ 2014; Pearson et al.\ 2018).

    \item[\textcolor{rojoquench}{Relación con aplanamiento}]
        La pendiente efectiva disminuye a $\alpha \approx 0.3$--$0.5$ para
        $M_\star > 10^{10}\,M_\odot$ a bajo $z$. Encontrada en estudios con SFR
        del IR o radio (Whitaker et al.\ 2014; Schreiber et al.\ 2015;
        Leslie et al.\ 2020).
\end{description}

El aplanamiento es más prominente a \textbf{bajo redshift} y puede explicarse
por la mezcla de galaxias morfológicamente distintas en la muestra.

\subsection{El modelo de Leslie et al.\ (2020)}

Se propone la siguiente parametrización que permite el aplanamiento:

\begin{equation}
    \langle\log\,\mathrm{SFR}\rangle =
    S_0(t) - \log\!\left[1 + \left(\frac{10^M}{10^{M'_t}}\right)^{-1}\right],
    \label{eq:ms_leslie}
\end{equation}

donde $M = \log(M_\star/M_\odot)$, $t$ es la edad del universo en Gyr, y:

\begin{align}
    S_0(t) &= S_0^{(1)} + a_1\,t, \label{eq:s0}\\
    M'_t   &= M_0 - a_2\,t.       \label{eq:mt}
\end{align}

Esta forma tiene dos regímenes límite naturales:

\begin{equation}
    \langle\log\,\mathrm{SFR}\rangle \approx
    \begin{cases}
        M + (S_0 - M'_t) & \text{si } M \ll M'_t
        \quad\text{(pendiente unitaria, sSFR constante)}\\[6pt]
        S_0(t)           & \text{si } M \gg M'_t
        \quad\text{(aplanamiento)}
    \end{cases}
    \label{eq:limites}
\end{equation}

\begin{center}
\begin{tikzpicture}
\begin{axis}[
    width=11cm, height=6.5cm,
    xlabel={$\log_{10}(M_\star\,/\,M_\odot)$},
    ylabel={$\log_{10}(\mathrm{SFR}\,/\,M_\odot\,\mathrm{yr}^{-1})$},
    xmin=8.5, xmax=12.0,
    ymin=-0.5, ymax=3.0,
    grid=both, grid style={gray!20},
    legend pos=north west,
    legend style={font=\small},
    title={\textbf{Forma de la MS según Leslie et al.\ (2020)}},
    title style={color=azulms}
]

%% z~0.4 (t~9.5 Gyr): S0 = -2.97 + 0.22*9.5 = -0.88, Mt = 11.06 - 0.12*9.5 = 9.92
\addplot[thick, color=rojoquench, domain=8.5:12.0, samples=150]
    { -0.88 - ln(1 + 10^(9.92 - x))/ln(10) };
\addlegendentry{$z \sim 0.4$ ($t \approx 9.5$\,Gyr)}

%% z~1.3 (t~4.7 Gyr): S0 = -2.97 + 0.22*4.7 = -1.93, Mt = 11.06 - 0.12*4.7 = 10.50
\addplot[thick, color=naranjabulge, domain=8.5:12.0, samples=150]
    { -1.93 - ln(1 + 10^(10.50 - x))/ln(10) };
\addlegendentry{$z \sim 1.3$ ($t \approx 4.7$\,Gyr)}

%% z~2.0 (t~3.3 Gyr): S0 = -2.97 + 0.22*3.3 = -2.24, Mt = 11.06 - 0.12*3.3 = 10.66
\addplot[thick, color=verdedisco, domain=8.5:12.0, samples=150]
    { -2.24 - ln(1 + 10^(10.66 - x))/ln(10) };
\addlegendentry{$z \sim 2.0$ ($t \approx 3.3$\,Gyr)}

%% z~3.0 (t~2.1 Gyr): S0 = -2.97 + 0.22*2.1 = -2.51, Mt = 11.06 - 0.12*2.1 = 10.81
\addplot[thick, color=azulms, domain=8.5:12.0, samples=150]
    { -2.51 - ln(1 + 10^(10.81 - x))/ln(10) };
\addlegendentry{$z \sim 3.0$ ($t \approx 2.1$\,Gyr)}

%% Flecha indicando evolución
\draw[->, thick, gray]
    (axis cs: 9.5, 0.6) -- (axis cs: 9.5, 1.9)
    node[midway, right, font=\small, gray] {$z$ crece};

\end{axis}
\end{tikzpicture}
\end{center}

Nótese cómo la MS se \textbf{desplaza hacia arriba} (mayor SFR) al aumentar
el redshift, y cómo el \textbf{aplanamiento} es más prominente a bajo $z$.

% ════════════════════════════════════════════════════════════════════════════
\section{Evolución con el Redshift: Valores Típicos de SFR}
% ════════════════════════════════════════════════════════════════════════════

La tabla siguiente muestra los SFR típicos predichos por el modelo de
Leslie et al.\ (2020) para galaxias SF con $M_\star = 10^{10}\,M_\odot$:

\begin{center}
\begin{tabular}{cccc}
\toprule
\textbf{Redshift} $z$ & \textbf{Edad} $t(z)$ & \textbf{SFR típico} & \textbf{sSFR} \\
& [Gyr] & [$M_\odot\,\mathrm{yr}^{-1}$] & [$\mathrm{Gyr}^{-1}$] \\
\midrule
0.4 & 9.5 & $\sim$1--2   & $\sim$0.1--0.2 \\
1.0 & 5.9 & $\sim$5--10  & $\sim$0.5--1.0 \\
2.0 & 3.3 & $\sim$30--50 & $\sim$3--5 \\
3.0 & 2.1 & $\sim$100    & $\sim$10 \\
5.0 & 1.2 & $\sim$300    & $\sim$30 \\
\bottomrule
\end{tabular}
\end{center}

La evolución es dramática: a $z \sim 3$ una galaxia de la misma masa
formaba estrellas $\sim 50$ veces más rápido que hoy.

% ════════════════════════════════════════════════════════════════════════════
\section{Morfología y la Secuencia Principal}
% ════════════════════════════════════════════════════════════════════════════

\subsection{Tipos morfológicos y su posición en la MS}

Leslie et al.\ (2020) muestran con claridad que la morfología de una galaxia
determina su posición relativa en la MS:

\begin{center}
\begin{tikzpicture}
\begin{axis}[
    width=11cm, height=6cm,
    xlabel={$\log_{10}(M_\star\,/\,M_\odot)$},
    ylabel={$\Delta\log_{10}(\mathrm{SFR})$ [dex]},
    xmin=9.5, xmax=11.5,
    ymin=-1.0, ymax=0.8,
    grid=both, grid style={gray!20},
    legend pos=south west,
    legend style={font=\small},
    title={\textbf{Desviación de la MS según morfología (z $\sim$ 0.75)}},
    title style={color=azulms},
    yticklabel style={/pgf/number format/fixed}
]

%% Línea de referencia MS
\addplot[thick, black, dashed, domain=9.5:11.5] {0};
\addlegendentry{MS promedio (referencia)}

%% Disco puro: MS más empinada → por encima de la media a altas masas
\addplot[thick, color=verdedisco, domain=9.5:11.2, samples=80]
    {0.25*(x - 10.5)};
\addlegendentry{Disco puro (más empinada)}

%% Dominadas por bulbo
\addplot[thick, color=naranjabulge, dashed, domain=9.5:11.5, samples=80]
    {-0.15*(x - 10.0) - 0.05};
\addlegendentry{Bulbo dominante}

%% Irregulares
\addplot[thick, color=moradostarburst, dotted, domain=9.5:11.0, samples=80]
    {0.35};
\addlegendentry{Irregulares (mergers)}

%% Early-type
\addplot[thick, color=rojoquench, domain=10.0:11.5, samples=80]
    {-0.5 - 0.1*(x - 10.5)};
\addlegendentry{Early-type (elípticas)}

\end{axis}
\end{tikzpicture}
\end{center}

\subsection{¿Por qué se aplana la MS a altas masas?}

El argumento central del paper es el siguiente:

\begin{enumerate}[leftmargin=2em]
    \item Las galaxias de \textbf{disco puro} siguen una MS \textit{lineal}
          (pendiente $\approx 0.75$) sin aplanamiento.
    \item Las galaxias \textbf{dominadas por bulbo} tienen menor SFR a masa
          fija porque el bulbo es quiescente (no forma estrellas).
    \item A bajo redshift, la fracción de galaxias con bulbo prominente
          \textbf{aumenta}, especialmente a altas masas.
    \item La \textit{combinación} de discos + bulbos produce el aplanamiento
          observado en la MS promedio.
\end{enumerate}

Matemáticamente, si el bulbo tiene fracción $B/T$ de la masa total y SFR $\approx 0$:

\begin{equation}
    \mathrm{SFR}_{\mathrm{total}} \approx \mathrm{SFR}_{\mathrm{disco}}
    = \mathrm{SFR}_{\mathrm{MS,disco}} \times (1 - B/T)^{\beta},
    \label{eq:bulge}
\end{equation}

donde $\beta$ depende de la pendiente de la MS de los discos. A mayor $B/T$,
menor SFR total a masa fija.

% ════════════════════════════════════════════════════════════════════════════
\section{Quenching: ¿Cómo se apagan las Galaxias?}
% ════════════════════════════════════════════════════════════════════════════

El ``apagado'' de la formación estelar (\textit{quenching}) es el proceso
por el cual una galaxia transita desde la MS hacia la región quiescente.
Se proponen varios mecanismos:

\begin{center}
\begin{tabular}{p{3.5cm} p{4cm} p{4.5cm}}
\toprule
\textbf{Mecanismo} & \textbf{Escala de masa} & \textbf{Descripción} \\
\midrule
\textit{Mass quenching} (Peng et al.\ 2010)
    & $M_\star \gtrsim 10^{10.2}\,M_\odot$
    & Feedback de AGN, calentamiento gravitacional \\[6pt]
\textit{Environmental quenching}
    & $M_\star \lesssim 10^{10}\,M_\odot$
    & Ram-pressure stripping, strangulation en cúmulos \\[6pt]
\textit{Morphological quenching}
    & $M_\star \gtrsim 10^{10}\,M_\odot$
    & El bulbo estabiliza el gas contra la fragmentación \\[6pt]
\textit{Halo quenching}
    & $M_h \gtrsim 10^{12}\,M_\odot$
    & Shock-heating impide acreción de gas frío \\
\bottomrule
\end{tabular}
\end{center}

Leslie et al.\ (2020) concluyen que el \textbf{quenching de masa} está ligado
a cambios morfológicos: la presencia del bulbo es suficiente para explicar la
reducción del SFR en galaxias masivas, sin necesidad de suprimir el SFR del disco.

% ════════════════════════════════════════════════════════════════════════════
\section{Downsizing Cósmico}
% ════════════════════════════════════════════════════════════════════════════

Un resultado clave que emerge de la MS y su evolución es el
\textbf{downsizing} o ``encogimiento cósmico'': las galaxias más masivas
formaron sus estrellas \textit{primero} y se apagaron antes.

Evidencias en el paper:
\begin{itemize}[leftmargin=2em]
    \item La \textbf{masa característica} $M_\star^{\rm char}$ que maximiza la
          contribución a la densidad de SFR cósmica \textit{aumenta} con el
          redshift (más masas grandes contribuyen más a alto $z$).
    \item La \textbf{masa de quiebre} $M'_t$ de la MS también aumenta con $z$:
          \begin{equation}
              M'_t = M_0 - a_2\,t \approx 11.06 - 0.12\,t,
          \end{equation}
          lo que significa que a alto $z$ el aplanamiento ocurre en masas
          mayores que a bajo $z$.
    \item El sSFR de las galaxias masivas disminuye más rápido con el tiempo
          que el de las galaxias de baja masa.
\end{itemize}

\begin{center}
\begin{tabular}{ccc}
\toprule
\textbf{Redshift} $z$ & \textbf{Masa característica} & \textbf{Masa de quiebre} $M'_t$ \\
& $\log(M_\star^{\rm char}/M_\odot)$ & $\log(M_\star/M_\odot)$ \\
\midrule
$\sim$0.35 & 9.9  & $\sim$10.0 \\
$\sim$1.0  & 10.2 & $\sim$10.4 \\
$\sim$2.0  & 10.5 & $\sim$10.7 \\
$\sim$2.5  & 10.6 & $\sim$10.8 \\
\bottomrule
\end{tabular}
\end{center}

% ════════════════════════════════════════════════════════════════════════════
\section{Resumen Conceptual}
% ════════════════════════════════════════════════════════════════════════════

\begin{center}
\colorbox{grisclaro}{%
\begin{minipage}{0.9\textwidth}
\vspace{0.5em}
\textbf{La Secuencia Principal en pocas ideas:}
\vspace{0.3em}
\begin{enumerate}[leftmargin=2em]
    \item La MS es una \textbf{correlación universal} entre SFR y $M_\star$
          para galaxias star-forming, con dispersión $\sigma \approx 0.2$--$0.3$\,dex.
    \item La \textbf{normalización aumenta con el redshift}: el universo
          joven formaba estrellas mucho más activamente.
    \item A \textbf{altas masas} y bajo $z$, la MS se \textbf{aplana}, reflejo
          de la creciente fracción de galaxias con bulbo dominante.
    \item El \textbf{sSFR} disminuye con la masa y aumenta con $z$,
          con $\mathrm{sSFR} \propto (1+z)^{2.8}$.
    \item Las galaxias más masivas se apagan antes (\textbf{downsizing}).
    \item El \textbf{entorno} modifica la fracción de quiescentes pero
          no la MS de las galaxias activas.
\end{enumerate}
\vspace{0.5em}
\end{minipage}}
\end{center}

\vspace{1cm}
\noindent\rule{\linewidth}{0.5pt}\\[0.3em]
{\small \textbf{Referencias principales:}
Noeske et al.\ (2007) ApJL 660, L43 $\cdot$
Karim et al.\ (2011) ApJ 730, 61 $\cdot$
Speagle et al.\ (2014) ApJS 214, 15 $\cdot$
Peng et al.\ (2010) ApJ 721, 193 $\cdot$
Leslie et al.\ (2020) ApJ 899, 58.}

\end{document}